\documentclass[11pt]{article}
\usepackage[margin=0.85in]{geometry}
\usepackage[T1]{fontenc}
\usepackage{lmodern,amsmath,amssymb,booktabs,array,enumitem,xcolor,hyperref}
\definecolor{elvia}{HTML}{264E69}
\definecolor{promptink}{HTML}{545F6A}
\hypersetup{colorlinks=true,urlcolor=elvia,linkcolor=elvia}
\setlength{\parindent}{0pt}
\setlength{\emergencystretch}{1em}
\setlength{\parskip}{4pt}
\setlist{nosep,leftmargin=*}
\newcommand{\prompt}[1]{\textit{\textcolor{promptink}{[#1]}}}
\newcommand{\entry}[2]{\textbf{#1}\quad\prompt{#2}\par}
\newcommand{\guidance}[1]{\begin{quote}\small #1\end{quote}}
\begin{document}
{\Large\bfseries\color{elvia} ELVIA: First paper and report}\par
{\small Editable scientific template, version 2026-10-23. This is a template,
not a completed review or an endorsement.}
\guidance{Replace bracketed prompts with your own report. A focused check of one
claim is sufficient; no minimum page count is required. Keep each problem numbered
and use the headings below. Delete inapplicable optional sections. Compile with
pdfLaTeX or XeLaTeX, or upload this single UTF-8 .tex document. Include all equations
and references here; do not depend on external images, input files or bibliography
files. If external material is essential to reading the report, submit a
self-contained PDF instead.}

\textbf{Submission files}
\begin{itemize}
\item This report: one PDF or one complete .tex file.
\item Original article: one PDF.
\item Main original supplementary document, if applicable: one separate PDF. Select Not applicable on the website when there is none.
\item Optional verification code: one ZIP or 1--10 source-code/support files,
at most 5 MB total, including unpacked ZIP contents. A short README with commands
and dependencies is helpful; no prescribed folder structure is required.
\end{itemize}
The report, original article, original supplement and verification code have
separate upload slots. Do not combine them into a report ZIP.

\section{Paper identity and scope}
\entry{Paper}{Exact title, authors, venue/year, DOI or official URL.}
\entry{Version}{Published version, arXiv version number or precise dated version.}
\entry{Materials inspected}{Main PDF version; supplement title/version, or none;
public author-code repository and exact commit/version, if inspected.}
\entry{Scope}{Which claims, sections, experiments or methods you checked and why.}
\entry{Summary}{A short account of what your check establishes and what remains
unresolved. Distinguish clarification, potential issue and demonstrated issue.}
\entry{What was not checked}{Unexamined experiments, assumptions, datasets or claims.}

\newpage
\section{Problem 1: \prompt{short descriptive title}}
\guidance{Repeat this entire section as Problem 2, Problem 3, etc.\ for additional
independent issues. Do not split one issue into many fragments. A well-supported
No issue found within scope result is also useful.}

\subsection{Location and assessment}
\entry{Exact location}{PDF page and section; equation/figure/table number; or code
file, function and lines plus repository commit.}
\entry{Severity/outcome}{Choose exactly one: Minor; Major; Needs evidence;
No issue found within scope.}
\begin{itemize}
\item \textbf{Minor}: a local error or missing detail with bounded impact;
do not claim the main conclusion fails without showing that consequence.
\item \textbf{Major}: demonstrated effect on a central claim or essential result.
Show the dependence between the issue and that claim.
\item \textbf{Needs evidence}: a plausible concern whose validity or impact
has not yet been established.
\item \textbf{No issue found within scope}: the specified check supports the
original argument. This is not a certificate for the whole paper.
\end{itemize}

\subsection{Claim and assumptions}
\entry{Original claim}{Brief exact quotation or faithful paraphrase, with location.}
\entry{Conditions}{Definitions, assumptions, units, parameter range, sample or
experimental unit, and which conditions are stated versus inferred.}

\subsection{Check and evidence}
\entry{Method}{Step-by-step derivation, counterexample, calculation, experimental
control, statistical check or code inspection. Explain why it tests the claim.}
\entry{Result}{Actual observations or derived results. Separate analytic proof,
finite numerical checks, empirical observations and conjecture.}
\guidance{For calculations, give the formula, symbols and substitutions, for example
\(r=x/x_0\), with units and the condition \(x_0\ne 0\). For data analyses, report
sample counts, exclusions, independent units, preprocessing and uncertainty.
For a computational check, give versions, commands, random seeds and tolerances where relevant,
expected output and what you actually obtained. Do not imply a program was run
when it was only inspected.}
\entry{Evidence identifiers}{References, code filenames, output values or stable
artifact links needed to inspect this check.}

\subsection{Impact, counterevidence and possible correction}
\entry{Impact}{Exactly which result changes, under what conditions, and which
claims remain unaffected. Justify the chosen severity.}
\entry{Counterevidence}{The strongest alternative reading or contrary evidence
you checked, and whether it resolves the concern.}
\entry{Limits}{Unresolved steps, unavailable information, failed checks and
what this evidence does not establish.}
\entry{Possible correction (optional)}{A precise clarification, corrected statement
or repair candidate. State what still needs verification; do not claim adoption.}

\section{Reproduction and overall conclusion}
\entry{Verification materials}{List your optional uploaded code files or ZIP and
what each does. If no code is needed, state that explicitly.}
\entry{Reproduction}{Commands, dependencies and expected output, or Not applicable
for a self-contained mathematical/textual check. Do not include credentials.}
\entry{Conclusion}{A concise, scoped conclusion. Separate issue assessment,
verification of a proposed repair, and whether authors have adopted any change.}

\section{AI assistance and contributor declaration}
\guidance{AI-assisted scientific review is encouraged. Use AI to organize the report,
check a derivation or inspect code, then verify the source, assumptions and outputs
yourself. Do not treat model agreement as independent evidence. Manual work is
equally welcome.}
\entry{AI used}{Yes or no. If no, delete the AI-specific entries below.}
\entry{Model and access}{Actual provider, model/version, interface or tool,
reasoning setting, and date(s) used. Do not invent unavailable settings.}
\entry{Assistance}{Which tasks AI performed; what text, sources or code it saw;
material prompts or a concise reproducible description.}
\entry{Human checks}{Which source passages, calculations, citations and outputs
you personally verified. Identify unverified AI suggestions explicitly.}
\entry{Conflicts and role}{Authorship, collaboration, financial or other relevant
relationships; your role in this analysis.}
\entry{Responsibility}{I checked the cited evidence and accurately described the
scope, uncertainty, AI assistance and limitations of this contribution.}
\entry{Attribution}{Contributor name or preferred credit, affiliation if relevant,
and date. Do not put private contact details in the report.}

\guidance{List only sources actually consulted. Give authors, title, year, DOI or
stable URL and exact location when relevant. Brief quotations must be clearly
marked. Use this self-contained bibliography; no separate .bib file is needed.}
\begin{thebibliography}{9}
\bibitem{paper} \prompt{Authors. Title. Venue, year. DOI or stable URL and version.}
\bibitem{other} \prompt{Additional source actually used. Delete if not applicable.}
\end{thebibliography}
\end{document}
